%%%%%%%%%%%%%%%%%%%%%%%%%%%%%%%%%%%%%%%%%%%%%%%%%%%%%%%%%%%%
% 14.12 TRANSFORM TABLES
% The Composition of Advanced Mathematics
%%%%%%%%%%%%%%%%%%%%%%%%%%%%%%%%%%%%%%%%%%%%%%%%%%%%%%%%%%%%

\section{Transform Tables}
\label{sec:transform-tables}

\prerequisite{
Integration, differential equations, complex exponentials,
series, piecewise functions, improper integrals, linear algebra,
and basic complex-variable notation.
}

\learningobjectives{
By the end of this reference section, the reader should be able to:
\begin{itemize}
    \item recognize integral transforms as changes of representation;
    \item use common Laplace-transform pairs;
    \item recognize inverse Laplace transforms;
    \item apply Laplace linearity;
    \item use time shifting and frequency shifting;
    \item use derivative and integral properties of Laplace transforms;
    \item use convolution;
    \item transform ordinary differential equations;
    \item recognize unilateral Laplace transforms and initial conditions;
    \item recognize Fourier-transform conventions;
    \item use common Fourier-transform pairs;
    \item use Fourier shift, scaling, modulation, and differentiation laws;
    \item recognize the convolution theorem;
    \item understand Parseval and Plancherel relationships;
    \item distinguish angular-frequency and ordinary-frequency conventions;
    \item recognize Fourier sine and cosine transforms;
    \item recognize the discrete Fourier transform;
    \item recognize the \(z\)-transform and common transform pairs;
    \item relate the \(z\)-transform to difference equations;
    \item recognize transform methods used in PDEs, control theory,
          signals, probability, and numerical mathematics;
    \item avoid common convention and scaling errors.
\end{itemize}
}

%%%%%%%%%%%%%%%%%%%%%%%%%%%%%%%%%%%%%%%%%%%%%%%%%%%%%%%%%%%%
\subsection{What Is an Integral Transform?}
\label{subsec:integral-transform-reference}
%%%%%%%%%%%%%%%%%%%%%%%%%%%%%%%%%%%%%%%%%%%%%%%%%%%%%%%%%%%%

An integral transform converts a function \(f\) into another function
\(F\) by an expression of form

\[
\boxed{
F(s)
=
\int
K(s,t)f(t)\,dt,
}
\]

where \(K\) is called the transform kernel.

Transforms often convert difficult operations into simpler ones.

For example:

\[
\boxed{
\text{differentiation}
\longrightarrow
\text{multiplication},
}
\]

or

\[
\boxed{
\text{convolution}
\longrightarrow
\text{multiplication}.
}
\]

%%%%%%%%%%%%%%%%%%%%%%%%%%%%%%%%%%%%%%%%%%%%%%%%%%%%%%%%%%%%
\subsection{Laplace Transform Definition}
\label{subsec:laplace-transform-definition-reference}
%%%%%%%%%%%%%%%%%%%%%%%%%%%%%%%%%%%%%%%%%%%%%%%%%%%%%%%%%%%%

For a suitable function \(f(t)\) defined for \(t\geq0\), the Laplace transform
is

\[
\boxed{
\mathcal{L}\{f(t)\}
=
F(s)
=
\int_0^\infty
e^{-st}f(t)\,dt.
}
\]

%%%%%%%%%%%%%%%%%%%%%%%%%%%%%%%%%%%%%%%%%%%%%%%%%%%%%%%%%%%%
\subsection{Inverse Laplace Transform}
\label{subsec:inverse-laplace-reference}
%%%%%%%%%%%%%%%%%%%%%%%%%%%%%%%%%%%%%%%%%%%%%%%%%%%%%%%%%%%%

If

\[
F(s)
=
\mathcal{L}\{f(t)\},
\]

then

\[
\boxed{
f(t)
=
\mathcal{L}^{-1}\{F(s)\}.
}
\]

In practical computations, inverse transforms are often obtained from tables,
partial fractions, convolution, or complex inversion methods.

%%%%%%%%%%%%%%%%%%%%%%%%%%%%%%%%%%%%%%%%%%%%%%%%%%%%%%%%%%%%
\subsection{Laplace Linearity}
\label{subsec:laplace-linearity-reference}
%%%%%%%%%%%%%%%%%%%%%%%%%%%%%%%%%%%%%%%%%%%%%%%%%%%%%%%%%%%%

For constants \(a,b\),

\[
\boxed{
\mathcal{L}
\{
af+bg
\}
=
a\mathcal{L}\{f\}
+
b\mathcal{L}\{g\}.
}
\]

%%%%%%%%%%%%%%%%%%%%%%%%%%%%%%%%%%%%%%%%%%%%%%%%%%%%%%%%%%%%
\subsection{Laplace Transform of a Constant}
\label{subsec:laplace-constant}
%%%%%%%%%%%%%%%%%%%%%%%%%%%%%%%%%%%%%%%%%%%%%%%%%%%%%%%%%%%%

For \(s>0\),

\[
\boxed{
\mathcal{L}\{1\}
=
\frac1s.
}
\]

More generally,

\[
\boxed{
\mathcal{L}\{c\}
=
\frac{c}{s}.
}
\]

%%%%%%%%%%%%%%%%%%%%%%%%%%%%%%%%%%%%%%%%%%%%%%%%%%%%%%%%%%%%
\subsection{Laplace Transform of Powers}
\label{subsec:laplace-power-reference}
%%%%%%%%%%%%%%%%%%%%%%%%%%%%%%%%%%%%%%%%%%%%%%%%%%%%%%%%%%%%

For nonnegative integer \(n\),

\[
\boxed{
\mathcal{L}\{t^n\}
=
\frac{
n!
}{
s^{n+1}
}.
}
\]

More generally, for suitable \(\alpha>-1\),

\[
\boxed{
\mathcal{L}\{t^\alpha\}
=
\frac{
\Gamma(\alpha+1)
}{
s^{\alpha+1}
}.
}
\]

%%%%%%%%%%%%%%%%%%%%%%%%%%%%%%%%%%%%%%%%%%%%%%%%%%%%%%%%%%%%
\subsection{Laplace Transform of an Exponential}
\label{subsec:laplace-exponential}
%%%%%%%%%%%%%%%%%%%%%%%%%%%%%%%%%%%%%%%%%%%%%%%%%%%%%%%%%%%%

\[
\boxed{
\mathcal{L}\{e^{at}\}
=
\frac{
1
}{
s-a
}
}
\]

for

\[
\operatorname{Re}(s)>\operatorname{Re}(a).
\]

%%%%%%%%%%%%%%%%%%%%%%%%%%%%%%%%%%%%%%%%%%%%%%%%%%%%%%%%%%%%
\subsection{Laplace Transform of Sine}
\label{subsec:laplace-sine}
%%%%%%%%%%%%%%%%%%%%%%%%%%%%%%%%%%%%%%%%%%%%%%%%%%%%%%%%%%%%

\[
\boxed{
\mathcal{L}\{\sin bt\}
=
\frac{
b
}{
s^2+b^2
}.
}
\]

%%%%%%%%%%%%%%%%%%%%%%%%%%%%%%%%%%%%%%%%%%%%%%%%%%%%%%%%%%%%
\subsection{Laplace Transform of Cosine}
\label{subsec:laplace-cosine}
%%%%%%%%%%%%%%%%%%%%%%%%%%%%%%%%%%%%%%%%%%%%%%%%%%%%%%%%%%%%

\[
\boxed{
\mathcal{L}\{\cos bt\}
=
\frac{
s
}{
s^2+b^2
}.
}
\]

%%%%%%%%%%%%%%%%%%%%%%%%%%%%%%%%%%%%%%%%%%%%%%%%%%%%%%%%%%%%
\subsection{Laplace Transform of Hyperbolic Sine}
\label{subsec:laplace-sinh}
%%%%%%%%%%%%%%%%%%%%%%%%%%%%%%%%%%%%%%%%%%%%%%%%%%%%%%%%%%%%

\[
\boxed{
\mathcal{L}\{\sinh bt\}
=
\frac{
b
}{
s^2-b^2
}.
}
\]

%%%%%%%%%%%%%%%%%%%%%%%%%%%%%%%%%%%%%%%%%%%%%%%%%%%%%%%%%%%%
\subsection{Laplace Transform of Hyperbolic Cosine}
\label{subsec:laplace-cosh}
%%%%%%%%%%%%%%%%%%%%%%%%%%%%%%%%%%%%%%%%%%%%%%%%%%%%%%%%%%%%

\[
\boxed{
\mathcal{L}\{\cosh bt\}
=
\frac{
s
}{
s^2-b^2
}.
}
\]

%%%%%%%%%%%%%%%%%%%%%%%%%%%%%%%%%%%%%%%%%%%%%%%%%%%%%%%%%%%%
\subsection{Quick Laplace Transform Table}
\label{subsec:quick-laplace-table}
%%%%%%%%%%%%%%%%%%%%%%%%%%%%%%%%%%%%%%%%%%%%%%%%%%%%%%%%%%%%

\[
\boxed{
\mathcal{L}\{1\}
=
\frac1s
}
\]

\[
\boxed{
\mathcal{L}\{t\}
=
\frac1{s^2}
}
\]

\[
\boxed{
\mathcal{L}\{t^n\}
=
\frac{n!}{s^{n+1}}
}
\]

\[
\boxed{
\mathcal{L}\{e^{at}\}
=
\frac1{s-a}
}
\]

\[
\boxed{
\mathcal{L}\{\sin bt\}
=
\frac{b}{s^2+b^2}
}
\]

\[
\boxed{
\mathcal{L}\{\cos bt\}
=
\frac{s}{s^2+b^2}
}
\]

\[
\boxed{
\mathcal{L}\{\sinh bt\}
=
\frac{b}{s^2-b^2}
}
\]

\[
\boxed{
\mathcal{L}\{\cosh bt\}
=
\frac{s}{s^2-b^2}.
}
\]

%%%%%%%%%%%%%%%%%%%%%%%%%%%%%%%%%%%%%%%%%%%%%%%%%%%%%%%%%%%%
\subsection{First Derivative Property}
\label{subsec:laplace-first-derivative}
%%%%%%%%%%%%%%%%%%%%%%%%%%%%%%%%%%%%%%%%%%%%%%%%%%%%%%%%%%%%

If

\[
F(s)=\mathcal{L}\{f(t)\},
\]

then

\[
\boxed{
\mathcal{L}\{f'(t)\}
=
sF(s)-f(0).
}
\]

%%%%%%%%%%%%%%%%%%%%%%%%%%%%%%%%%%%%%%%%%%%%%%%%%%%%%%%%%%%%
\subsection{Second Derivative Property}
\label{subsec:laplace-second-derivative}
%%%%%%%%%%%%%%%%%%%%%%%%%%%%%%%%%%%%%%%%%%%%%%%%%%%%%%%%%%%%

\[
\boxed{
\mathcal{L}\{f''(t)\}
=
s^2F(s)
-
sf(0)
-
f'(0).
}
\]

%%%%%%%%%%%%%%%%%%%%%%%%%%%%%%%%%%%%%%%%%%%%%%%%%%%%%%%%%%%%
\subsection{Higher Laplace Derivatives}
\label{subsec:laplace-higher-derivatives}
%%%%%%%%%%%%%%%%%%%%%%%%%%%%%%%%%%%%%%%%%%%%%%%%%%%%%%%%%%%%

For positive integer \(n\),

\[
\boxed{
\mathcal{L}\{f^{(n)}(t)\}
=
s^nF(s)
-
s^{n-1}f(0)
-
s^{n-2}f'(0)
-\cdots
-
f^{(n-1)}(0).
}
\]

This property makes the Laplace transform especially useful for initial-value
problems.

%%%%%%%%%%%%%%%%%%%%%%%%%%%%%%%%%%%%%%%%%%%%%%%%%%%%%%%%%%%%
\subsection{Integral Property}
\label{subsec:laplace-integral-property}
%%%%%%%%%%%%%%%%%%%%%%%%%%%%%%%%%%%%%%%%%%%%%%%%%%%%%%%%%%%%

If

\[
g(t)
=
\int_0^t
f(\tau)\,d\tau,
\]

then

\[
\boxed{
\mathcal{L}\{g(t)\}
=
\frac{
F(s)
}{
s
}.
}
\]

%%%%%%%%%%%%%%%%%%%%%%%%%%%%%%%%%%%%%%%%%%%%%%%%%%%%%%%%%%%%
\subsection{Multiplication by \(t\)}
\label{subsec:laplace-multiplication-by-t}
%%%%%%%%%%%%%%%%%%%%%%%%%%%%%%%%%%%%%%%%%%%%%%%%%%%%%%%%%%%%

If

\[
F(s)
=
\mathcal{L}\{f(t)\},
\]

then

\[
\boxed{
\mathcal{L}\{t f(t)\}
=
-\frac{dF}{ds}.
}
\]

More generally,

\[
\boxed{
\mathcal{L}\{t^nf(t)\}
=
(-1)^n
\frac{
d^nF
}{
ds^n
}.
}
\]

%%%%%%%%%%%%%%%%%%%%%%%%%%%%%%%%%%%%%%%%%%%%%%%%%%%%%%%%%%%%
\subsection{First Shifting Theorem}
\label{subsec:laplace-first-shift}
%%%%%%%%%%%%%%%%%%%%%%%%%%%%%%%%%%%%%%%%%%%%%%%%%%%%%%%%%%%%

If

\[
F(s)
=
\mathcal{L}\{f(t)\},
\]

then

\[
\boxed{
\mathcal{L}\{e^{at}f(t)\}
=
F(s-a).
}
\]

%%%%%%%%%%%%%%%%%%%%%%%%%%%%%%%%%%%%%%%%%%%%%%%%%%%%%%%%%%%%
\subsection{Unit Step Function}
\label{subsec:unit-step-reference}
%%%%%%%%%%%%%%%%%%%%%%%%%%%%%%%%%%%%%%%%%%%%%%%%%%%%%%%%%%%%

The Heaviside unit step is

\[
\boxed{
u(t-a)
=
\begin{cases}
0, & t<a,\\
1, & t>a.
\end{cases}
}
\]

The value at \(t=a\) is convention dependent and usually irrelevant for
ordinary transform calculations.

%%%%%%%%%%%%%%%%%%%%%%%%%%%%%%%%%%%%%%%%%%%%%%%%%%%%%%%%%%%%
\subsection{Second Shifting Theorem}
\label{subsec:laplace-second-shift}
%%%%%%%%%%%%%%%%%%%%%%%%%%%%%%%%%%%%%%%%%%%%%%%%%%%%%%%%%%%%

If

\[
F(s)
=
\mathcal{L}\{f(t)\},
\]

then

\[
\boxed{
\mathcal{L}
\{
u(t-a)f(t-a)
\}
=
e^{-as}F(s).
}
\]

%%%%%%%%%%%%%%%%%%%%%%%%%%%%%%%%%%%%%%%%%%%%%%%%%%%%%%%%%%%%
\subsection{Impulse Transform}
\label{subsec:dirac-delta-laplace-reference}
%%%%%%%%%%%%%%%%%%%%%%%%%%%%%%%%%%%%%%%%%%%%%%%%%%%%%%%%%%%%

For the Dirac delta in the generalized-function sense,

\[
\boxed{
\mathcal{L}\{\delta(t-a)\}
=
e^{-as},
\qquad
a\geq0.
}
\]

In particular,

\[
\boxed{
\mathcal{L}\{\delta(t)\}
=
1.
}
\]

%%%%%%%%%%%%%%%%%%%%%%%%%%%%%%%%%%%%%%%%%%%%%%%%%%%%%%%%%%%%
\subsection{Laplace Convolution}
\label{subsec:laplace-convolution-reference}
%%%%%%%%%%%%%%%%%%%%%%%%%%%%%%%%%%%%%%%%%%%%%%%%%%%%%%%%%%%%

The convolution of \(f\) and \(g\) on \([0,\infty)\) is

\[
\boxed{
(f*g)(t)
=
\int_0^t
f(\tau)
g(t-\tau)
\,d\tau.
}
\]

%%%%%%%%%%%%%%%%%%%%%%%%%%%%%%%%%%%%%%%%%%%%%%%%%%%%%%%%%%%%
\subsection{Laplace Convolution Theorem}
\label{subsec:laplace-convolution-theorem}
%%%%%%%%%%%%%%%%%%%%%%%%%%%%%%%%%%%%%%%%%%%%%%%%%%%%%%%%%%%%

If

\[
F(s)=\mathcal{L}\{f\},
\qquad
G(s)=\mathcal{L}\{g\},
\]

then

\[
\boxed{
\mathcal{L}\{f*g\}
=
F(s)G(s).
}
\]

Therefore,

\[
\boxed{
\mathcal{L}^{-1}\{F(s)G(s)\}
=
f*g.
}
\]

%%%%%%%%%%%%%%%%%%%%%%%%%%%%%%%%%%%%%%%%%%%%%%%%%%%%%%%%%%%%
\subsection{Periodic-Function Laplace Transform}
\label{subsec:laplace-periodic-function}
%%%%%%%%%%%%%%%%%%%%%%%%%%%%%%%%%%%%%%%%%%%%%%%%%%%%%%%%%%%%

If \(f\) has period \(T>0\), then under suitable conditions,

\[
\boxed{
\mathcal{L}\{f\}
=
\frac{
\displaystyle
\int_0^T
e^{-st}f(t)\,dt
}{
1-e^{-sT}
}.
}
\]

%%%%%%%%%%%%%%%%%%%%%%%%%%%%%%%%%%%%%%%%%%%%%%%%%%%%%%%%%%%%
\subsection{Worked Example: Laplace Transform of an ODE}
\label{subsec:worked-laplace-ode}
%%%%%%%%%%%%%%%%%%%%%%%%%%%%%%%%%%%%%%%%%%%%%%%%%%%%%%%%%%%%

\begin{workedexamplebox}{Solve a Simple Initial-Value Problem}

Solve

\[
y'+y=0,
\qquad
y(0)=1.
\]

Take the Laplace transform:

\[
sY(s)-1+Y(s)=0.
\]

Thus

\[
(s+1)Y(s)=1.
\]

Therefore,

\[
Y(s)
=
\frac{1}{s+1}.
\]

Taking the inverse transform,

\begin{answerbox}
\[
\boxed{
y(t)=e^{-t}.
}
\]
\end{answerbox}

\end{workedexamplebox}

%%%%%%%%%%%%%%%%%%%%%%%%%%%%%%%%%%%%%%%%%%%%%%%%%%%%%%%%%%%%
\subsection{Inverse Laplace Through Partial Fractions}
\label{subsec:inverse-laplace-partial-fractions}
%%%%%%%%%%%%%%%%%%%%%%%%%%%%%%%%%%%%%%%%%%%%%%%%%%%%%%%%%%%%

Rational transforms are often inverted by decomposing them into standard
pieces.

For example,

\[
\frac{1}{s(s+1)}
=
\frac1s
-
\frac1{s+1}.
\]

Thus

\[
\boxed{
\mathcal{L}^{-1}
\left\{
\frac{1}{s(s+1)}
\right\}
=
1-e^{-t}.
}
\]

%%%%%%%%%%%%%%%%%%%%%%%%%%%%%%%%%%%%%%%%%%%%%%%%%%%%%%%%%%%%
\subsection{Fourier Transform Convention}
\label{subsec:fourier-transform-convention-reference}
%%%%%%%%%%%%%%%%%%%%%%%%%%%%%%%%%%%%%%%%%%%%%%%%%%%%%%%%%%%%

One common angular-frequency convention is

\[
\boxed{
\widehat{f}(\omega)
=
\int_{-\infty}^{\infty}
f(x)e^{-i\omega x}
\,dx.
}
\]

The inverse transform is

\[
\boxed{
f(x)
=
\frac{1}{2\pi}
\int_{-\infty}^{\infty}
\widehat{f}(\omega)
e^{i\omega x}
\,d\omega.
}
\]

Fourier-transform conventions vary between texts.

%%%%%%%%%%%%%%%%%%%%%%%%%%%%%%%%%%%%%%%%%%%%%%%%%%%%%%%%%%%%
\subsection{Symmetric Fourier Convention}
\label{subsec:symmetric-fourier-convention}
%%%%%%%%%%%%%%%%%%%%%%%%%%%%%%%%%%%%%%%%%%%%%%%%%%%%%%%%%%%%

Another common convention uses

\[
\boxed{
\widehat{f}(\omega)
=
\frac{1}{\sqrt{2\pi}}
\int_{-\infty}^{\infty}
f(x)e^{-i\omega x}
\,dx,
}
\]

with inverse

\[
\boxed{
f(x)
=
\frac{1}{\sqrt{2\pi}}
\int_{-\infty}^{\infty}
\widehat{f}(\omega)
e^{i\omega x}
\,d\omega.
}
\]

Always check the convention before using a transform table.

%%%%%%%%%%%%%%%%%%%%%%%%%%%%%%%%%%%%%%%%%%%%%%%%%%%%%%%%%%%%
\subsection{Ordinary-Frequency Fourier Convention}
\label{subsec:ordinary-frequency-fourier}
%%%%%%%%%%%%%%%%%%%%%%%%%%%%%%%%%%%%%%%%%%%%%%%%%%%%%%%%%%%%

Using frequency \(\nu\) rather than angular frequency \(\omega\),

\[
\boxed{
\widehat{f}(\nu)
=
\int_{-\infty}^{\infty}
f(t)e^{-2\pi i\nu t}
\,dt.
}
\]

Then

\[
\boxed{
f(t)
=
\int_{-\infty}^{\infty}
\widehat{f}(\nu)
e^{2\pi i\nu t}
\,d\nu.
}
\]

%%%%%%%%%%%%%%%%%%%%%%%%%%%%%%%%%%%%%%%%%%%%%%%%%%%%%%%%%%%%
\subsection{Fourier Linearity}
\label{subsec:fourier-linearity-reference}
%%%%%%%%%%%%%%%%%%%%%%%%%%%%%%%%%%%%%%%%%%%%%%%%%%%%%%%%%%%%

\[
\boxed{
\mathcal{F}\{af+bg\}
=
a\widehat f
+
b\widehat g.
}
\]

%%%%%%%%%%%%%%%%%%%%%%%%%%%%%%%%%%%%%%%%%%%%%%%%%%%%%%%%%%%%
\subsection{Fourier Transform of a Gaussian}
\label{subsec:fourier-gaussian-reference}
%%%%%%%%%%%%%%%%%%%%%%%%%%%%%%%%%%%%%%%%%%%%%%%%%%%%%%%%%%%%

Under the angular-frequency convention

\[
\widehat f(\omega)
=
\int_{-\infty}^{\infty}
f(x)e^{-i\omega x}
\,dx,
\]

one has

\[
\boxed{
\mathcal{F}\{e^{-ax^2}\}
=
\sqrt{
\frac{\pi}{a}
}
e^{-\omega^2/(4a)}
}
\]

for \(a>0\).

Thus Gaussians transform into Gaussians.

%%%%%%%%%%%%%%%%%%%%%%%%%%%%%%%%%%%%%%%%%%%%%%%%%%%%%%%%%%%%
\subsection{Fourier Transform of a Decaying Exponential}
\label{subsec:fourier-decaying-exponential}
%%%%%%%%%%%%%%%%%%%%%%%%%%%%%%%%%%%%%%%%%%%%%%%%%%%%%%%%%%%%

For \(a>0\),

\[
f(x)
=
e^{-a|x|}
\]

has transform

\[
\boxed{
\widehat f(\omega)
=
\frac{
2a
}{
a^2+\omega^2
}.
}
\]

%%%%%%%%%%%%%%%%%%%%%%%%%%%%%%%%%%%%%%%%%%%%%%%%%%%%%%%%%%%%
\subsection{Fourier Transform of a Rectangular Pulse}
\label{subsec:fourier-rectangular-pulse}
%%%%%%%%%%%%%%%%%%%%%%%%%%%%%%%%%%%%%%%%%%%%%%%%%%%%%%%%%%%%

Let

\[
f(x)
=
\begin{cases}
1, & |x|\leq a,\\
0, & |x|>a.
\end{cases}
\]

Then

\[
\boxed{
\widehat f(\omega)
=
\frac{
2\sin(a\omega)
}{
\omega
}
}
\]

with the value at \(\omega=0\) interpreted by continuity as \(2a\).

%%%%%%%%%%%%%%%%%%%%%%%%%%%%%%%%%%%%%%%%%%%%%%%%%%%%%%%%%%%%
\subsection{Fourier Time Shift}
\label{subsec:fourier-time-shift}
%%%%%%%%%%%%%%%%%%%%%%%%%%%%%%%%%%%%%%%%%%%%%%%%%%%%%%%%%%%%

If

\[
g(x)=f(x-a),
\]

then

\[
\boxed{
\widehat g(\omega)
=
e^{-i\omega a}
\widehat f(\omega).
}
\]

Translation in physical space becomes multiplication by a phase factor in
frequency space.

%%%%%%%%%%%%%%%%%%%%%%%%%%%%%%%%%%%%%%%%%%%%%%%%%%%%%%%%%%%%
\subsection{Fourier Frequency Shift}
\label{subsec:fourier-frequency-shift}
%%%%%%%%%%%%%%%%%%%%%%%%%%%%%%%%%%%%%%%%%%%%%%%%%%%%%%%%%%%%

If

\[
g(x)
=
e^{iax}f(x),
\]

then

\[
\boxed{
\widehat g(\omega)
=
\widehat f(\omega-a).
}
\]

This is also called modulation.

%%%%%%%%%%%%%%%%%%%%%%%%%%%%%%%%%%%%%%%%%%%%%%%%%%%%%%%%%%%%
\subsection{Fourier Scaling}
\label{subsec:fourier-scaling}
%%%%%%%%%%%%%%%%%%%%%%%%%%%%%%%%%%%%%%%%%%%%%%%%%%%%%%%%%%%%

For nonzero real \(a\),

\[
\boxed{
\mathcal{F}
\{
f(ax)
\}
(\omega)
=
\frac1{|a|}
\widehat f
\left(
\frac{\omega}{a}
\right).
}
\]

%%%%%%%%%%%%%%%%%%%%%%%%%%%%%%%%%%%%%%%%%%%%%%%%%%%%%%%%%%%%
\subsection{Fourier Derivative Property}
\label{subsec:fourier-derivative-property}
%%%%%%%%%%%%%%%%%%%%%%%%%%%%%%%%%%%%%%%%%%%%%%%%%%%%%%%%%%%%

Under suitable conditions,

\[
\boxed{
\mathcal{F}\{f'(x)\}
=
i\omega
\widehat f(\omega).
}
\]

More generally,

\[
\boxed{
\mathcal{F}\{f^{(n)}(x)\}
=
(i\omega)^n
\widehat f(\omega).
}
\]

%%%%%%%%%%%%%%%%%%%%%%%%%%%%%%%%%%%%%%%%%%%%%%%%%%%%%%%%%%%%
\subsection{Multiplication by \(x\) in Fourier Space}
\label{subsec:fourier-multiplication-x}
%%%%%%%%%%%%%%%%%%%%%%%%%%%%%%%%%%%%%%%%%%%%%%%%%%%%%%%%%%%%

Under suitable regularity,

\[
\boxed{
\mathcal{F}\{x f(x)\}
=
i
\frac{d}{d\omega}
\widehat f(\omega).
}
\]

%%%%%%%%%%%%%%%%%%%%%%%%%%%%%%%%%%%%%%%%%%%%%%%%%%%%%%%%%%%%
\subsection{Fourier Convolution}
\label{subsec:fourier-convolution-reference}
%%%%%%%%%%%%%%%%%%%%%%%%%%%%%%%%%%%%%%%%%%%%%%%%%%%%%%%%%%%%

For functions on \(\mathbb{R}\),

\[
\boxed{
(f*g)(x)
=
\int_{-\infty}^{\infty}
f(t)
g(x-t)
\,dt.
}
\]

%%%%%%%%%%%%%%%%%%%%%%%%%%%%%%%%%%%%%%%%%%%%%%%%%%%%%%%%%%%%
\subsection{Fourier Convolution Theorem}
\label{subsec:fourier-convolution-theorem}
%%%%%%%%%%%%%%%%%%%%%%%%%%%%%%%%%%%%%%%%%%%%%%%%%%%%%%%%%%%%

Under the convention used here,

\[
\boxed{
\mathcal{F}\{f*g\}
=
\widehat f(\omega)
\widehat g(\omega).
}
\]

Also,

\[
\boxed{
\mathcal{F}\{fg\}
=
\frac{1}{2\pi}
(\widehat f*\widehat g).
}
\]

The constants depend on the Fourier-transform convention.

%%%%%%%%%%%%%%%%%%%%%%%%%%%%%%%%%%%%%%%%%%%%%%%%%%%%%%%%%%%%
\subsection{Parseval Formula}
\label{subsec:fourier-parseval-transform}
%%%%%%%%%%%%%%%%%%%%%%%%%%%%%%%%%%%%%%%%%%%%%%%%%%%%%%%%%%%%

Under suitable hypotheses,

\[
\boxed{
\int_{-\infty}^{\infty}
f(x)
\overline{g(x)}
\,dx
=
\frac{1}{2\pi}
\int_{-\infty}^{\infty}
\widehat f(\omega)
\overline{
\widehat g(\omega)
}
\,d\omega.
}
\]

%%%%%%%%%%%%%%%%%%%%%%%%%%%%%%%%%%%%%%%%%%%%%%%%%%%%%%%%%%%%
\subsection{Plancherel Formula}
\label{subsec:plancherel-transform-reference}
%%%%%%%%%%%%%%%%%%%%%%%%%%%%%%%%%%%%%%%%%%%%%%%%%%%%%%%%%%%%

Taking \(g=f\),

\[
\boxed{
\int_{-\infty}^{\infty}
|f(x)|^2
\,dx
=
\frac{1}{2\pi}
\int_{-\infty}^{\infty}
|\widehat f(\omega)|^2
\,d\omega.
}
\]

Thus the Fourier transform preserves energy up to normalization.

%%%%%%%%%%%%%%%%%%%%%%%%%%%%%%%%%%%%%%%%%%%%%%%%%%%%%%%%%%%%
\subsection{Fourier Transform of the Delta Distribution}
\label{subsec:fourier-delta-reference}
%%%%%%%%%%%%%%%%%%%%%%%%%%%%%%%%%%%%%%%%%%%%%%%%%%%%%%%%%%%%

In the generalized-function sense,

\[
\boxed{
\mathcal{F}\{\delta(x)\}
=
1.
}
\]

Also,

\[
\boxed{
\mathcal{F}\{\delta(x-a)\}
=
e^{-i\omega a}.
}
\]

%%%%%%%%%%%%%%%%%%%%%%%%%%%%%%%%%%%%%%%%%%%%%%%%%%%%%%%%%%%%
\subsection{Fourier Transform of a Constant}
\label{subsec:fourier-constant-distribution}
%%%%%%%%%%%%%%%%%%%%%%%%%%%%%%%%%%%%%%%%%%%%%%%%%%%%%%%%%%%%

In the distributional sense,

\[
\boxed{
\mathcal{F}\{1\}
=
2\pi\delta(\omega)
}
\]

under the angular-frequency convention used here.

%%%%%%%%%%%%%%%%%%%%%%%%%%%%%%%%%%%%%%%%%%%%%%%%%%%%%%%%%%%%
\subsection{Fourier Transform of Complex Exponentials}
\label{subsec:fourier-complex-exponential}
%%%%%%%%%%%%%%%%%%%%%%%%%%%%%%%%%%%%%%%%%%%%%%%%%%%%%%%%%%%%

Distributionally,

\[
\boxed{
\mathcal{F}
\{
e^{iax}
\}
=
2\pi
\delta(\omega-a).
}
\]

This expresses a pure sinusoid as a single frequency.

%%%%%%%%%%%%%%%%%%%%%%%%%%%%%%%%%%%%%%%%%%%%%%%%%%%%%%%%%%%%
\subsection{Sine and Cosine in Frequency Space}
\label{subsec:fourier-sine-cosine-delta}
%%%%%%%%%%%%%%%%%%%%%%%%%%%%%%%%%%%%%%%%%%%%%%%%%%%%%%%%%%%%

Using

\[
\cos ax
=
\frac12
(e^{iax}+e^{-iax}),
\]

\[
\boxed{
\mathcal{F}\{\cos ax\}
=
\pi
[
\delta(\omega-a)
+
\delta(\omega+a)
].
}
\]

Similarly,

\[
\boxed{
\mathcal{F}\{\sin ax\}
=
\frac{\pi}{i}
[
\delta(\omega-a)
-
\delta(\omega+a)
].
}
\]

%%%%%%%%%%%%%%%%%%%%%%%%%%%%%%%%%%%%%%%%%%%%%%%%%%%%%%%%%%%%
\subsection{Fourier Sine Transform}
\label{subsec:fourier-sine-transform-reference}
%%%%%%%%%%%%%%%%%%%%%%%%%%%%%%%%%%%%%%%%%%%%%%%%%%%%%%%%%%%%

One common convention is

\[
\boxed{
F_s(\omega)
=
\sqrt{
\frac{2}{\pi}
}
\int_0^\infty
f(x)\sin(\omega x)
\,dx.
}
\]

It is especially useful for odd extensions and boundary-value problems.

%%%%%%%%%%%%%%%%%%%%%%%%%%%%%%%%%%%%%%%%%%%%%%%%%%%%%%%%%%%%
\subsection{Fourier Cosine Transform}
\label{subsec:fourier-cosine-transform-reference}
%%%%%%%%%%%%%%%%%%%%%%%%%%%%%%%%%%%%%%%%%%%%%%%%%%%%%%%%%%%%

One common convention is

\[
\boxed{
F_c(\omega)
=
\sqrt{
\frac{2}{\pi}
}
\int_0^\infty
f(x)\cos(\omega x)
\,dx.
}
\]

It is especially useful for even extensions.

%%%%%%%%%%%%%%%%%%%%%%%%%%%%%%%%%%%%%%%%%%%%%%%%%%%%%%%%%%%%
\subsection{Discrete Fourier Transform}
\label{subsec:discrete-fourier-transform-reference}
%%%%%%%%%%%%%%%%%%%%%%%%%%%%%%%%%%%%%%%%%%%%%%%%%%%%%%%%%%%%

For data

\[
x_0,x_1,\ldots,x_{N-1},
\]

the discrete Fourier transform is

\[
\boxed{
X_k
=
\sum_{n=0}^{N-1}
x_n
e^{-2\pi i kn/N},
\qquad
k=0,\ldots,N-1.
}
\]

%%%%%%%%%%%%%%%%%%%%%%%%%%%%%%%%%%%%%%%%%%%%%%%%%%%%%%%%%%%%
\subsection{Inverse Discrete Fourier Transform}
\label{subsec:inverse-dft-reference}
%%%%%%%%%%%%%%%%%%%%%%%%%%%%%%%%%%%%%%%%%%%%%%%%%%%%%%%%%%%%

\[
\boxed{
x_n
=
\frac1N
\sum_{k=0}^{N-1}
X_k
e^{2\pi i kn/N}.
}
\]

%%%%%%%%%%%%%%%%%%%%%%%%%%%%%%%%%%%%%%%%%%%%%%%%%%%%%%%%%%%%
\subsection{Fast Fourier Transform}
\label{subsec:fft-reference}
%%%%%%%%%%%%%%%%%%%%%%%%%%%%%%%%%%%%%%%%%%%%%%%%%%%%%%%%%%%%

The Fast Fourier Transform, or FFT, is an algorithm for computing the DFT
efficiently.

A direct DFT requires approximately

\[
\boxed{
O(N^2)
}
\]

arithmetic operations.

FFT algorithms reduce this to approximately

\[
\boxed{
O(N\log N).
}
\]

%%%%%%%%%%%%%%%%%%%%%%%%%%%%%%%%%%%%%%%%%%%%%%%%%%%%%%%%%%%%
\subsection{Discrete Convolution Theorem}
\label{subsec:dft-convolution-theorem}
%%%%%%%%%%%%%%%%%%%%%%%%%%%%%%%%%%%%%%%%%%%%%%%%%%%%%%%%%%%%

Circular convolution in the discrete domain corresponds to multiplication of
DFTs.

If

\[
z=x*y
\]

denotes circular convolution, then

\[
\boxed{
Z_k
=
X_kY_k.
}
\]

%%%%%%%%%%%%%%%%%%%%%%%%%%%%%%%%%%%%%%%%%%%%%%%%%%%%%%%%%%%%
\subsection{\(z\)-Transform Definition}
\label{subsec:z-transform-definition-reference}
%%%%%%%%%%%%%%%%%%%%%%%%%%%%%%%%%%%%%%%%%%%%%%%%%%%%%%%%%%%%

For a sequence

\[
x[n],
\]

the bilateral \(z\)-transform is

\[
\boxed{
X(z)
=
\sum_{n=-\infty}^{\infty}
x[n]z^{-n}.
}
\]

%%%%%%%%%%%%%%%%%%%%%%%%%%%%%%%%%%%%%%%%%%%%%%%%%%%%%%%%%%%%
\subsection{Unilateral \(z\)-Transform}
\label{subsec:unilateral-z-transform}
%%%%%%%%%%%%%%%%%%%%%%%%%%%%%%%%%%%%%%%%%%%%%%%%%%%%%%%%%%%%

For sequences indexed by \(n\geq0\),

\[
\boxed{
X(z)
=
\sum_{n=0}^{\infty}
x[n]z^{-n}.
}
\]

This form is useful for initial-value difference equations.

%%%%%%%%%%%%%%%%%%%%%%%%%%%%%%%%%%%%%%%%%%%%%%%%%%%%%%%%%%%%
\subsection{\(z\)-Transform of the Unit Impulse}
\label{subsec:z-transform-unit-impulse}
%%%%%%%%%%%%%%%%%%%%%%%%%%%%%%%%%%%%%%%%%%%%%%%%%%%%%%%%%%%%

For the discrete impulse

\[
\delta[n],
\]

\[
\boxed{
\mathcal{Z}\{\delta[n]\}
=
1.
}
\]

%%%%%%%%%%%%%%%%%%%%%%%%%%%%%%%%%%%%%%%%%%%%%%%%%%%%%%%%%%%%
\subsection{\(z\)-Transform of the Unit Step}
\label{subsec:z-transform-unit-step}
%%%%%%%%%%%%%%%%%%%%%%%%%%%%%%%%%%%%%%%%%%%%%%%%%%%%%%%%%%%%

For

\[
u[n]
=
\begin{cases}
1, & n\geq0,\\
0, & n<0,
\end{cases}
\]

\[
\boxed{
\mathcal{Z}\{u[n]\}
=
\frac{1}{1-z^{-1}}
=
\frac{z}{z-1}
}
\]

with region of convergence

\[
\boxed{
|z|>1.
}
\]

%%%%%%%%%%%%%%%%%%%%%%%%%%%%%%%%%%%%%%%%%%%%%%%%%%%%%%%%%%%%
\subsection{\(z\)-Transform of a Geometric Sequence}
\label{subsec:z-transform-geometric}
%%%%%%%%%%%%%%%%%%%%%%%%%%%%%%%%%%%%%%%%%%%%%%%%%%%%%%%%%%%%

For

\[
x[n]
=
a^n u[n],
\]

\[
\boxed{
X(z)
=
\frac{1}{1-az^{-1}}
=
\frac{z}{z-a}
}
\]

with region of convergence

\[
\boxed{
|z|>|a|.
}
\]

%%%%%%%%%%%%%%%%%%%%%%%%%%%%%%%%%%%%%%%%%%%%%%%%%%%%%%%%%%%%
\subsection{\(z\)-Transform Linearity}
\label{subsec:z-transform-linearity}
%%%%%%%%%%%%%%%%%%%%%%%%%%%%%%%%%%%%%%%%%%%%%%%%%%%%%%%%%%%%

\[
\boxed{
\mathcal{Z}\{ax[n]+by[n]\}
=
aX(z)+bY(z).
}
\]

%%%%%%%%%%%%%%%%%%%%%%%%%%%%%%%%%%%%%%%%%%%%%%%%%%%%%%%%%%%%
\subsection{\(z\)-Transform Time Shift}
\label{subsec:z-transform-time-shift}
%%%%%%%%%%%%%%%%%%%%%%%%%%%%%%%%%%%%%%%%%%%%%%%%%%%%%%%%%%%%

For the bilateral transform,

\[
\boxed{
\mathcal{Z}
\{
x[n-n_0]
\}
=
z^{-n_0}X(z).
}
\]

%%%%%%%%%%%%%%%%%%%%%%%%%%%%%%%%%%%%%%%%%%%%%%%%%%%%%%%%%%%%
\subsection{\(z\)-Transform Convolution}
\label{subsec:z-transform-convolution}
%%%%%%%%%%%%%%%%%%%%%%%%%%%%%%%%%%%%%%%%%%%%%%%%%%%%%%%%%%%%

For discrete convolution

\[
(x*y)[n]
=
\sum_{k=-\infty}^{\infty}
x[k]y[n-k],
\]

\[
\boxed{
\mathcal{Z}\{x*y\}
=
X(z)Y(z).
}
\]

%%%%%%%%%%%%%%%%%%%%%%%%%%%%%%%%%%%%%%%%%%%%%%%%%%%%%%%%%%%%
\subsection{Difference Equations and the \(z\)-Transform}
\label{subsec:z-transform-difference-equations}
%%%%%%%%%%%%%%%%%%%%%%%%%%%%%%%%%%%%%%%%%%%%%%%%%%%%%%%%%%%%

The \(z\)-transform converts shifts such as

\[
x[n-1]
\]

into multiplication by powers of \(z^{-1}\).

Therefore linear difference equations can often be transformed into algebraic
equations in \(X(z)\).

%%%%%%%%%%%%%%%%%%%%%%%%%%%%%%%%%%%%%%%%%%%%%%%%%%%%%%%%%%%%
\subsection{Worked Example: Simple \(z\)-Transform}
\label{subsec:worked-z-transform}
%%%%%%%%%%%%%%%%%%%%%%%%%%%%%%%%%%%%%%%%%%%%%%%%%%%%%%%%%%%%

\begin{workedexamplebox}{Transform a Geometric Sequence}

Let

\[
x[n]
=
\left(
\frac12
\right)^n
u[n].
\]

Then

\[
X(z)
=
\sum_{n=0}^{\infty}
\left(
\frac12
\right)^n
z^{-n}.
\]

This is geometric with ratio

\[
\frac{1}{2z}.
\]

Therefore,

\begin{answerbox}
\[
\boxed{
X(z)
=
\frac{
1
}{
1-\frac12z^{-1}
}
=
\frac{
z
}{
z-\frac12
},
\qquad
|z|>\frac12.
}
\]
\end{answerbox}

\end{workedexamplebox}

%%%%%%%%%%%%%%%%%%%%%%%%%%%%%%%%%%%%%%%%%%%%%%%%%%%%%%%%%%%%
\subsection{Laplace Versus Fourier Transform}
\label{subsec:laplace-versus-fourier-reference}
%%%%%%%%%%%%%%%%%%%%%%%%%%%%%%%%%%%%%%%%%%%%%%%%%%%%%%%%%%%%

The Laplace transform uses the kernel

\[
\boxed{
e^{-st}
}
\]

with complex variable

\[
s=\sigma+i\omega.
\]

The Fourier transform corresponds formally to the imaginary axis

\[
\boxed{
s=i\omega
}
\]

when the relevant convergence conditions hold.

Thus the Laplace transform includes exponential weighting not present in the
ordinary Fourier transform.

%%%%%%%%%%%%%%%%%%%%%%%%%%%%%%%%%%%%%%%%%%%%%%%%%%%%%%%%%%%%
\subsection{Fourier Versus \(z\)-Transform}
\label{subsec:fourier-versus-z-transform}
%%%%%%%%%%%%%%%%%%%%%%%%%%%%%%%%%%%%%%%%%%%%%%%%%%%%%%%%%%%%

The \(z\)-transform is the discrete-time analogue of the Laplace transform.

For

\[
z=e^{i\omega},
\]

evaluation of the \(z\)-transform on the unit circle is related to the
discrete-time Fourier transform, when the region of convergence includes the
unit circle.

%%%%%%%%%%%%%%%%%%%%%%%%%%%%%%%%%%%%%%%%%%%%%%%%%%%%%%%%%%%%
\subsection{Transforms and Differential Equations}
\label{subsec:transforms-differential-equations}
%%%%%%%%%%%%%%%%%%%%%%%%%%%%%%%%%%%%%%%%%%%%%%%%%%%%%%%%%%%%

The Laplace transform converts

\[
\boxed{
\text{time derivatives}
}
\]

into algebraic expressions in \(s\).

The Fourier transform converts spatial derivatives into powers of
\(i\omega\).

This allows many differential equations to become algebraic equations in
transform space.

%%%%%%%%%%%%%%%%%%%%%%%%%%%%%%%%%%%%%%%%%%%%%%%%%%%%%%%%%%%%
\subsection{Fourier Transform and the Heat Equation}
\label{subsec:fourier-heat-equation-reference}
%%%%%%%%%%%%%%%%%%%%%%%%%%%%%%%%%%%%%%%%%%%%%%%%%%%%%%%%%%%%

For

\[
u_t
=
\kappa u_{xx},
\]

taking the Fourier transform in \(x\) gives

\[
\boxed{
\widehat{u}_t
=
-\kappa\omega^2
\widehat{u}.
}
\]

This is an ordinary differential equation in \(t\).

Its solution is

\[
\boxed{
\widehat u(\omega,t)
=
\widehat u(\omega,0)
e^{-\kappa\omega^2t}.
}
\]

%%%%%%%%%%%%%%%%%%%%%%%%%%%%%%%%%%%%%%%%%%%%%%%%%%%%%%%%%%%%
\subsection{Fourier Transform and the Wave Equation}
\label{subsec:fourier-wave-equation-reference}
%%%%%%%%%%%%%%%%%%%%%%%%%%%%%%%%%%%%%%%%%%%%%%%%%%%%%%%%%%%%

For

\[
u_{tt}
=
c^2u_{xx},
\]

the spatial Fourier transform gives

\[
\boxed{
\widehat u_{tt}
+
c^2\omega^2
\widehat u
=
0.
}
\]

Each Fourier frequency behaves as a harmonic oscillator.

%%%%%%%%%%%%%%%%%%%%%%%%%%%%%%%%%%%%%%%%%%%%%%%%%%%%%%%%%%%%
\subsection{Laplace Transform and Transfer Functions}
\label{subsec:laplace-transfer-function-reference}
%%%%%%%%%%%%%%%%%%%%%%%%%%%%%%%%%%%%%%%%%%%%%%%%%%%%%%%%%%%%

For a linear time-invariant system with zero initial conditions,

\[
\boxed{
H(s)
=
\frac{
Y(s)
}{
U(s)
}
}
\]

is the transfer function.

Here \(U(s)\) is the input transform and \(Y(s)\) is the output transform.

%%%%%%%%%%%%%%%%%%%%%%%%%%%%%%%%%%%%%%%%%%%%%%%%%%%%%%%%%%%%
\subsection{Impulse Response and Transfer Function}
\label{subsec:impulse-response-transfer-function}
%%%%%%%%%%%%%%%%%%%%%%%%%%%%%%%%%%%%%%%%%%%%%%%%%%%%%%%%%%%%

If \(h(t)\) is the impulse response, then

\[
\boxed{
H(s)
=
\mathcal{L}\{h(t)\}.
}
\]

The output is convolution:

\[
\boxed{
y(t)
=
(h*u)(t).
}
\]

In transform space,

\[
\boxed{
Y(s)
=
H(s)U(s).
}
\]

%%%%%%%%%%%%%%%%%%%%%%%%%%%%%%%%%%%%%%%%%%%%%%%%%%%%%%%%%%%%
\subsection{Transforms in Probability}
\label{subsec:transforms-probability-reference}
%%%%%%%%%%%%%%%%%%%%%%%%%%%%%%%%%%%%%%%%%%%%%%%%%%%%%%%%%%%%

For random variable \(X\), the characteristic function is

\[
\boxed{
\varphi_X(t)
=
\mathbb{E}
[
e^{itX}
].
}
\]

It is a Fourier-type transform of the probability distribution.

%%%%%%%%%%%%%%%%%%%%%%%%%%%%%%%%%%%%%%%%%%%%%%%%%%%%%%%%%%%%
\subsection{Moment-Generating Function}
\label{subsec:moment-generating-transform-reference}
%%%%%%%%%%%%%%%%%%%%%%%%%%%%%%%%%%%%%%%%%%%%%%%%%%%%%%%%%%%%

The moment-generating function is

\[
\boxed{
M_X(t)
=
\mathbb{E}
[
e^{tX}
].
}
\]

When it exists near \(t=0\),

\[
\boxed{
M_X^{(n)}(0)
=
\mathbb{E}[X^n].
}
\]

It is closely related to a bilateral Laplace transform.

%%%%%%%%%%%%%%%%%%%%%%%%%%%%%%%%%%%%%%%%%%%%%%%%%%%%%%%%%%%%
\subsection{Characteristic Function of a Normal Variable}
\label{subsec:normal-characteristic-function}
%%%%%%%%%%%%%%%%%%%%%%%%%%%%%%%%%%%%%%%%%%%%%%%%%%%%%%%%%%%%

If

\[
X\sim N(\mu,\sigma^2),
\]

then

\[
\boxed{
\varphi_X(t)
=
\exp
\left(
i\mu t
-
\frac12\sigma^2t^2
\right).
}
\]

%%%%%%%%%%%%%%%%%%%%%%%%%%%%%%%%%%%%%%%%%%%%%%%%%%%%%%%%%%%%
\subsection{Quick Laplace Property Table}
\label{subsec:quick-laplace-properties}
%%%%%%%%%%%%%%%%%%%%%%%%%%%%%%%%%%%%%%%%%%%%%%%%%%%%%%%%%%%%

If

\[
F(s)=\mathcal{L}\{f(t)\},
\]

then

\[
\boxed{
\mathcal{L}\{f'(t)\}
=
sF(s)-f(0)
}
\]

\[
\boxed{
\mathcal{L}\{e^{at}f(t)\}
=
F(s-a)
}
\]

\[
\boxed{
\mathcal{L}\{t f(t)\}
=
-F'(s)
}
\]

\[
\boxed{
\mathcal{L}\{f*g\}
=
F(s)G(s)
}
\]

\[
\boxed{
\mathcal{L}
\{
u(t-a)f(t-a)
\}
=
e^{-as}F(s).
}
\]

%%%%%%%%%%%%%%%%%%%%%%%%%%%%%%%%%%%%%%%%%%%%%%%%%%%%%%%%%%%%
\subsection{Quick Fourier Property Table}
\label{subsec:quick-fourier-properties}
%%%%%%%%%%%%%%%%%%%%%%%%%%%%%%%%%%%%%%%%%%%%%%%%%%%%%%%%%%%%

Under

\[
\widehat f(\omega)
=
\int_{-\infty}^{\infty}
f(x)e^{-i\omega x}\,dx,
\]

\[
\boxed{
f(x-a)
\longleftrightarrow
e^{-i\omega a}\widehat f(\omega)
}
\]

\[
\boxed{
e^{iax}f(x)
\longleftrightarrow
\widehat f(\omega-a)
}
\]

\[
\boxed{
f(ax)
\longleftrightarrow
\frac1{|a|}
\widehat f\left(\frac{\omega}{a}\right)
}
\]

\[
\boxed{
f'(x)
\longleftrightarrow
i\omega\widehat f(\omega)
}
\]

\[
\boxed{
xf(x)
\longleftrightarrow
i\widehat f'(\omega)
}
\]

\[
\boxed{
f*g
\longleftrightarrow
\widehat f\,\widehat g.
}
\]

%%%%%%%%%%%%%%%%%%%%%%%%%%%%%%%%%%%%%%%%%%%%%%%%%%%%%%%%%%%%
\subsection{Common Errors}
\label{subsec:transform-tables-common-errors}
%%%%%%%%%%%%%%%%%%%%%%%%%%%%%%%%%%%%%%%%%%%%%%%%%%%%%%%%%%%%

\subsubsection{Mixing Fourier Conventions}

Different texts place \(2\pi\) factors in different locations.

Never combine formulas from different conventions without adjusting their
normalization.

%%%%%%%%%%%%%%%%%%%%%%%%%%%%%%%%%%%%%%%%%%%%%%%%%%%%%%%%%%%%
\subsubsection{Confusing Frequency With Angular Frequency}

Ordinary frequency \(\nu\) and angular frequency \(\omega\) satisfy

\[
\boxed{
\omega
=
2\pi\nu.
}
\]

This changes the exponential kernel and transform scaling.

%%%%%%%%%%%%%%%%%%%%%%%%%%%%%%%%%%%%%%%%%%%%%%%%%%%%%%%%%%%%
\subsubsection{Forgetting Initial Conditions in Laplace Derivatives}

The transform

\[
\mathcal{L}\{f'\}
\]

is not simply \(sF(s)\).

It is

\[
\boxed{
sF(s)-f(0).
}
\]

%%%%%%%%%%%%%%%%%%%%%%%%%%%%%%%%%%%%%%%%%%%%%%%%%%%%%%%%%%%%
\subsubsection{Incorrect Transform of a Shift}

A time shift and a frequency shift produce different effects.

For Fourier transforms,

\[
f(x-a)
\]

produces a phase factor, while

\[
e^{iax}f(x)
\]

shifts the frequency-domain function.

%%%%%%%%%%%%%%%%%%%%%%%%%%%%%%%%%%%%%%%%%%%%%%%%%%%%%%%%%%%%
\subsubsection{Ignoring the Region of Convergence}

For Laplace and \(z\)-transforms, the algebraic expression alone may not
uniquely identify the original function or sequence.

The region of convergence is part of the transform description.

%%%%%%%%%%%%%%%%%%%%%%%%%%%%%%%%%%%%%%%%%%%%%%%%%%%%%%%%%%%%
\subsubsection{Forgetting the Absolute Value in Fourier Scaling}

For real nonzero \(a\),

\[
\boxed{
f(ax)
\longleftrightarrow
\frac1{|a|}
\widehat f(\omega/a).
}
\]

%%%%%%%%%%%%%%%%%%%%%%%%%%%%%%%%%%%%%%%%%%%%%%%%%%%%%%%%%%%%
\subsubsection{Using the Convolution Theorem With the Wrong Normalization}

Depending on the Fourier convention, factors such as

\[
2\pi
\]

may appear in the convolution or product theorem.

%%%%%%%%%%%%%%%%%%%%%%%%%%%%%%%%%%%%%%%%%%%%%%%%%%%%%%%%%%%%
\subsubsection{Treating the Dirac Delta as an Ordinary Function}

The Dirac delta is a distribution or generalized function.

Formulas involving

\[
\delta(x)
\]

should be interpreted in the appropriate distributional sense.

%%%%%%%%%%%%%%%%%%%%%%%%%%%%%%%%%%%%%%%%%%%%%%%%%%%%%%%%%%%%
\subsubsection{Confusing Fourier Series and Fourier Transforms}

Fourier series describe periodic functions using discrete frequencies.

Fourier transforms typically describe nonperiodic functions using a
continuous frequency variable.

%%%%%%%%%%%%%%%%%%%%%%%%%%%%%%%%%%%%%%%%%%%%%%%%%%%%%%%%%%%%
\subsubsection{Confusing the DFT With the Continuous Fourier Transform}

The DFT operates on a finite sequence and yields a finite collection of
frequency coefficients.

It introduces periodicity and sampling effects not present in the continuous
transform.

%%%%%%%%%%%%%%%%%%%%%%%%%%%%%%%%%%%%%%%%%%%%%%%%%%%%%%%%%%%%
\subsubsection{Ignoring Aliasing}

When continuous signals are sampled, frequencies above the Nyquist limit can
be represented incorrectly as lower frequencies.

Sampling theory must be considered in numerical Fourier analysis.

%%%%%%%%%%%%%%%%%%%%%%%%%%%%%%%%%%%%%%%%%%%%%%%%%%%%%%%%%%%%
\subsubsection{Assuming Every Function Has an Ordinary Fourier Transform}

Some important transforms exist only in:

\[
\boxed{
L^2
}
\]

or distributional senses.

For example, constants and pure sinusoids naturally involve delta
distributions.

%%%%%%%%%%%%%%%%%%%%%%%%%%%%%%%%%%%%%%%%%%%%%%%%%%%%%%%%%%%%
\subsection{Transform Method Workflow}
\label{subsec:transform-method-workflow}
%%%%%%%%%%%%%%%%%%%%%%%%%%%%%%%%%%%%%%%%%%%%%%%%%%%%%%%%%%%%

A practical transform workflow is:

\begin{enumerate}
    \item identify the problem domain: continuous time, continuous space,
          discrete time, or periodic;
    \item choose Laplace, Fourier, Fourier series, DFT, or \(z\)-transform;
    \item state the transform convention;
    \item check convergence or regularity assumptions;
    \item apply linearity;
    \item use shifting, scaling, or derivative properties;
    \item simplify algebraically in transform space;
    \item use partial fractions when appropriate;
    \item apply an inverse transform;
    \item restore initial or boundary conditions;
    \item verify the solution in the original domain.
\end{enumerate}

%%%%%%%%%%%%%%%%%%%%%%%%%%%%%%%%%%%%%%%%%%%%%%%%%%%%%%%%%%%%
\subsection{Exercises}
\label{subsec:transform-tables-exercises}
%%%%%%%%%%%%%%%%%%%%%%%%%%%%%%%%%%%%%%%%%%%%%%%%%%%%%%%%%%%%

\begin{exercise}[1]
Find

\[
\mathcal{L}\{1\}.
\]
\end{exercise}

\begin{exercise}[1]
Find

\[
\mathcal{L}\{t^3\}.
\]
\end{exercise}

\begin{exercise}[1]
Find

\[
\mathcal{L}\{e^{2t}\}.
\]
\end{exercise}

\begin{exercise}[1]
Find

\[
\mathcal{L}\{\sin3t\}.
\]
\end{exercise}

\begin{exercise}[1]
Find

\[
\mathcal{L}\{\cos4t\}.
\]
\end{exercise}

\begin{exercise}[2]
Find

\[
\mathcal{L}^{-1}
\left\{
\frac1{s^3}
\right\}.
\]
\end{exercise}

\begin{exercise}[2]
Find

\[
\mathcal{L}^{-1}
\left\{
\frac{5}{s^2+25}
\right\}.
\]
\end{exercise}

\begin{exercise}[2]
Find

\[
\mathcal{L}^{-1}
\left\{
\frac{s}{s^2+9}
\right\}.
\]
\end{exercise}

\begin{exercise}[2]
Use the first shifting theorem to find

\[
\mathcal{L}\{e^{2t}\sin3t\}.
\]
\end{exercise}

\begin{exercise}[2]
Use

\[
\mathcal{L}\{tf(t)\}
=
-F'(s)
\]

to transform \(t\sin t\).
\end{exercise}

\begin{exercise}[2]
Find the transform of

\[
u(t-2).
\]
\end{exercise}

\begin{exercise}[2]
Find the transform of

\[
u(t-a)(t-a).
\]
\end{exercise}

\begin{exercise}[2]
Compute a simple Laplace convolution.
\end{exercise}

\begin{exercise}[3]
Solve

\[
y'+2y=0,
\qquad
y(0)=3
\]

using Laplace transforms.
\end{exercise}

\begin{exercise}[3]
Solve

\[
y''+y=0,
\qquad
y(0)=0,
\quad
y'(0)=1
\]

using Laplace transforms.
\end{exercise}

\begin{exercise}[3]
Invert

\[
\frac{1}{s(s+2)}
\]

using partial fractions.
\end{exercise}

\begin{exercise}[3]
Find the Laplace transform of a periodic square-wave function using the
periodic-function formula.
\end{exercise}

\begin{exercise}[2]
State the Fourier-transform convention used in this section.
\end{exercise}

\begin{exercise}[2]
Find the effect of replacing \(f(x)\) by \(f(x-a)\) on its Fourier transform.
\end{exercise}

\begin{exercise}[2]
Find the effect of multiplying \(f(x)\) by \(e^{iax}\).
\end{exercise}

\begin{exercise}[2]
Find the Fourier transform of \(f'(x)\) in terms of \(\widehat f\).
\end{exercise}

\begin{exercise}[2]
State the Fourier convolution theorem.
\end{exercise}

\begin{exercise}[3]
Derive the Fourier transform of

\[
e^{-a|x|},
\qquad
a>0.
\]
\end{exercise}

\begin{exercise}[3]
Derive the Fourier transform of a rectangular pulse.
\end{exercise}

\begin{exercise}[3]
Verify the scaling property of the Fourier transform.
\end{exercise}

\begin{exercise}[3]
Use the Fourier transform to reduce the heat equation to an ODE in time.
\end{exercise}

\begin{exercise}[3]
Use the Fourier transform to reduce the wave equation to a family of harmonic
oscillators.
\end{exercise}

\begin{exercise}[3]
Compute the DFT of the sequence

\[
1,0,0,0.
\]
\end{exercise}

\begin{exercise}[3]
Compute the DFT of a constant length-\(N\) sequence.
\end{exercise}

\begin{exercise}[3]
Verify the inverse DFT formula for a small example.
\end{exercise}

\begin{exercise}[3]
Explain why FFT algorithms are computationally important.
\end{exercise}

\begin{exercise}[2]
Find the \(z\)-transform of

\[
\delta[n].
\]
\end{exercise}

\begin{exercise}[2]
Find the \(z\)-transform of

\[
u[n].
\]
\end{exercise}

\begin{exercise}[2]
Find the \(z\)-transform of

\[
3^n u[n].
\]
\end{exercise}

\begin{exercise}[3]
Determine the region of convergence for

\[
a^n u[n].
\]
\end{exercise}

\begin{exercise}[3]
Use the \(z\)-transform to solve a simple first-order difference equation.
\end{exercise}

\begin{exercise}[3]
Show that convolution in discrete time becomes multiplication in the
\(z\)-domain.
\end{exercise}

\begin{exercise}[3]
Compare the Laplace transform with the Fourier transform by setting

\[
s=\sigma+i\omega.
\]
\end{exercise}

\begin{exercise}[3]
Explain how the discrete-time Fourier transform is related to the
\(z\)-transform on the unit circle.
\end{exercise}

\begin{exercise}[3]
Find the characteristic function of a normal random variable.
\end{exercise}

\begin{exercise}[3]
Explain how derivatives of a moment-generating function produce moments.
\end{exercise}

\begin{exercise}[4]
Derive the Laplace derivative property from integration by parts.
\end{exercise}

\begin{exercise}[4]
Prove the Laplace convolution theorem.
\end{exercise}

\begin{exercise}[4]
Prove the Fourier shift theorem.
\end{exercise}

\begin{exercise}[4]
Prove the Fourier scaling theorem.
\end{exercise}

\begin{exercise}[4]
Prove the Fourier convolution theorem under suitable hypotheses.
\end{exercise}

\begin{exercise}[4]
Derive Plancherel's identity for a suitable class of functions.
\end{exercise}

\begin{exercise}[4]
Study the distributional Fourier transform of constants and sinusoids.
\end{exercise}

\begin{exercise}[4]
Derive the Fourier transform of a Gaussian.
\end{exercise}

\begin{exercise}[4]
Study the relationship among continuous Fourier transforms, sampled signals,
the DFT, and aliasing.
\end{exercise}

%%%%%%%%%%%%%%%%%%%%%%%%%%%%%%%%%%%%%%%%%%%%%%%%%%%%%%%%%%%%
\subsection{Conceptual Summary}
\label{subsec:transform-tables-summary}
%%%%%%%%%%%%%%%%%%%%%%%%%%%%%%%%%%%%%%%%%%%%%%%%%%%%%%%%%%%%

The Laplace transform is

\[
\boxed{
\mathcal{L}\{f(t)\}
=
\int_0^\infty
e^{-st}f(t)\,dt.
}
\]

Its derivative property is

\[
\boxed{
\mathcal{L}\{f'\}
=
sF(s)-f(0),
}
\]

which makes initial-value differential equations algebraic in transform space.

The Fourier transform, under one common convention, is

\[
\boxed{
\widehat f(\omega)
=
\int_{-\infty}^{\infty}
f(x)e^{-i\omega x}
\,dx.
}
\]

It converts differentiation into multiplication:

\[
\boxed{
f'
\longleftrightarrow
i\omega\widehat f.
}
\]

Both Laplace and Fourier transforms convert convolution into multiplication.

The DFT provides a finite discrete-frequency representation:

\[
\boxed{
X_k
=
\sum_{n=0}^{N-1}
x_n e^{-2\pi i kn/N}.
}
\]

The \(z\)-transform is

\[
\boxed{
X(z)
=
\sum_n
x[n]z^{-n},
}
\]

and is fundamental for discrete-time systems and difference equations.

The most important practical principles are

\[
\boxed{
\text{state the transform convention},
}
\]

\[
\boxed{
\text{track convergence and regions of convergence},
}
\]

and

\[
\boxed{
\text{use transforms to replace differential or convolution operations by
algebraic ones}.
}
\]

%%%%%%%%%%%%%%%%%%%%%%%%%%%%%%%%%%%%%%%%%%%%%%%%%%%%%%%%%%%%
\subsection{Section Connections}
\label{subsec:transform-tables-connections}
%%%%%%%%%%%%%%%%%%%%%%%%%%%%%%%%%%%%%%%%%%%%%%%%%%%%%%%%%%%%

\chapterconnection{
Transform Tables consolidates transform methods used throughout Differential
Equations, Analysis, Numerical Mathematics, Probability, Statistics,
Mathematical Physics, Control Theory, and signal processing. Laplace
transforms convert initial-value differential equations into algebraic
equations; Fourier transforms diagonalize constant-coefficient differential
operators and separate frequency content; convolution theorems connect
transforms with linear systems; the DFT and FFT connect continuous Fourier
analysis with computation; and the \(z\)-transform provides the discrete-time
analogue used for difference equations and digital systems. The next
reference section, Probability Distributions, collects the major discrete
and continuous probability distributions together with their parameters,
supports, moments, and common formulas.
}

%%%%%%%%%%%%%%%%%%%%%%%%%%%%%%%%%%%%%%%%%%%%%%%%%%%%%%%%%%%%
% SECTION SOURCE TRACKING
%%%%%%%%%%%%%%%%%%%%%%%%%%%%%%%%%%%%%%%%%%%%%%%%%%%%%%%%%%%%

%------------------------------------------------------------
% 14.12 Transform Tables
%------------------------------------------------------------
%
% Source basis:
%   - Standard Laplace-transform theory.
%   - Standard Fourier analysis.
%   - Standard discrete Fourier and z-transform theory.
%
% Used for:
%   - integral transforms;
%   - Laplace transforms;
%   - inverse Laplace transforms;
%   - Laplace transform pairs;
%   - derivative properties;
%   - integral properties;
%   - first and second shifting theorems;
%   - unit-step functions;
%   - Dirac delta transforms;
%   - convolution;
%   - periodic functions;
%   - ODE solution by transforms;
%   - Fourier-transform conventions;
%   - Gaussian transforms;
%   - pulse transforms;
%   - shift, modulation, scaling, and differentiation;
%   - Fourier convolution;
%   - Parseval and Plancherel formulas;
%   - distributional transforms;
%   - sine and cosine transforms;
%   - discrete Fourier transforms;
%   - inverse DFT;
%   - FFT;
%   - discrete convolution;
%   - z-transforms;
%   - unilateral z-transforms;
%   - regions of convergence;
%   - difference equations;
%   - PDE transform methods;
%   - transfer functions;
%   - impulse responses;
%   - characteristic functions;
%   - moment-generating functions;
%   - common errors;
%   - applications and exercises.
%
% Notes:
%   - Probability Distributions is developed next in Section 14.13.
%   - Statistical Tables follows in Section 14.14.
%
%%%%%%%%%%%%%%%%%%%%%%%%%%%%%%%%%%%%%%%%%%%%%%%%%%%%%%%%%%%%